% ══════════════════════════════════════════════════════════════════════════════
% Lecture 10: Dynamic Programming – Longest Common Subsequence (LCS)
% ══════════════════════════════════════════════════════════════════════════════

% ══════════════════════════════════════════════════════════════════════════════
\section*{Longest Common Subsequence (LCS)}
\addcontentsline{toc}{subsection}{Longest Common Subsequence (LCS)}
% ══════════════════════════════════════════════════════════════════════════════

% ─────────────────────────────────────────────────────────────────────────────
\subsection{Definition and Motivation}

\begin{definition}{Longest Common Subsequence}{lcs}
\textbf{Input:} two sequences $X = \langle x_1, \ldots, x_m \rangle$ and
$Y = \langle y_1, \ldots, y_n \rangle$.

\textbf{Output:} a common subsequence to both whose length is maximal.

\textbf{Reminder:} a subsequence does not need to be consecutive, but must
respect the order of the elements.
\end{definition}

\textbf{Example:} for the words \texttt{heroically} and \texttt{scholarly}, the LCS
is \texttt{holly} (length~5): the letters appear in the correct order in
each of the two words, without necessarily being adjacent.

% ─────────────────────────────────────────────────────────────────────────────
\subsection{Why Naive Approaches Fail}

\textbf{Brute force:} for each subsequence of $X$, check if it is
a subsequence of $Y$.

\begin{complexity}{Brute force for LCS}{lcsbf}
\begin{itemize}[noitemsep]
  \item $X$ has $2^m$ subsequences.
  \item Each verification takes $\Theta(n)$ (traversing $Y$ letter by letter).
  \item \textbf{Total complexity: $\Theta(n \cdot 2^m)$ — exponential, unusable.}
\end{itemize}
\end{complexity}

There is also no natural greedy algorithm for this problem:
the locally best letter does not necessarily lead to the global optimal
solution. Dynamic programming is the right tool.

% ─────────────────────────────────────────────────────────────────────────────
\subsection{Intuition of the DP Approach}

We traverse the two sequences from the \emph{right} (the end) to the left,
step by step.

At each position $(i, j)$ (looking at $x_i$ and $y_j$):
\begin{itemize}
  \item \textbf{If $x_i = y_j$}: this letter is part of the LCS; we include it
        and solve the subproblem on prefixes $X_{i-1}$ and $Y_{j-1}$.
  \item \textbf{If $x_i \ne y_j$}: the optimal LCS is obtained by moving one step
        to the left in \emph{one} of the two sequences, and taking the best of the
        two results.
\end{itemize}

\begin{center}
\begin{tikzpicture}[font=\small, >=Stealth]
  % X sequence
  \foreach \c/\i in {B/0,A/1,B/2,D/3,B/4,A/5}{
    \node[draw, fill=blue!10, minimum width=0.65cm, minimum height=0.55cm]
          at (\i*0.72, 0.8) {\texttt{\c}};
  }
  \node[left=4pt] at (-0.4, 0.8) {$X$:};
  % Y sequence
  \foreach \c/\i in {D/0,A/1,C/2,B/3,C/4,B/5,A/6}{
    \node[draw, fill=orange!15, minimum width=0.65cm, minimum height=0.55cm]
          at (\i*0.72, 0) {\texttt{\c}};
  }
  \node[left=4pt] at (-0.4, 0) {$Y$:};
  % Arrows for matching letters
  \draw[->, keygreen, thick] (0*0.72, 0.52) -- (3*0.72, 0.28)
        node[midway, right=2pt, font=\tiny, keygreen]{B};
  \draw[->, keyred, thick, dashed] (5*0.72, 0.52) -- (1*0.72, 0.28)
        node[midway, left=10pt, font=\tiny, keyred]{A};
\end{tikzpicture}
\end{center}

% ─────────────────────────────────────────────────────────────────────────────
\subsection{Optimal Substructure}

Let $X_i = \langle x_1, x_2, \ldots, x_i \rangle$ and
$Y_j = \langle y_1, y_2, \ldots, y_j \rangle$ be the prefixes of length $i$ and $j$.

\begin{theorem}{Optimal Substructure for LCS}{lcsopt}
Let $Z = \langle z_1, z_2, \ldots, z_k \rangle$ be an LCS of $X_i$ and $Y_j$.
\begin{enumerate}
  \item \textbf{If $x_i = y_j$}: then $z_k = x_i = y_j$ and $Z_{k-1}$ is an
        LCS of $X_{i-1}$ and $Y_{j-1}$.
  \item \textbf{If $x_i \ne y_j$ and $z_k \ne x_i$}: then $Z$ is an LCS of
        $X_{i-1}$ and $Y_j$.
  \item \textbf{If $x_i \ne y_j$ and $z_k \ne y_j$}: then $Z$ is an LCS of
        $X_i$ and $Y_{j-1}$.
\end{enumerate}
\end{theorem}

\begin{proof}
\textbf{Case 1 ($x_i = y_j$).}
If $z_k \ne x_i = y_j$, then $Z' = \langle z_1, \ldots, z_k, x_i \rangle$ is a
common subsequence of $X_i$ and $Y_j$ of length $k+1$, which contradicts
the optimality of $Z$. Thus $z_k = x_i = y_j$.

Now suppose that $Z_{k-1}$ is not an LCS of $X_{i-1}$ and $Y_{j-1}$:
there would exist a common subsequence $W$ of $X_{i-1}$ and $Y_{j-1}$ of
length $\ge k$. Then $\langle W, z_k \rangle$ would be a common subsequence
of $X_i$ and $Y_j$ of length $\ge k+1$, a contradiction.

\textbf{Case 2 ($x_i \ne y_j$, $z_k \ne x_i$).}
$Z$ is already a common subsequence of $X_{i-1}$ and $Y_j$. If it were not
optimal, there would exist $W$ of length $> k$ as a subsequence of $X_{i-1}$ and $Y_j$,
which would also be a subsequence of $X_i$ and $Y_j$ (since $W$ does not end with $x_i$),
contradicting the optimality of $Z$. Case~3 is symmetric.
\end{proof}

% ─────────────────────────────────────────────────────────────────────────────
\subsection{Recursive Formulation}

It follows from the previous theorem that the length of the LCS of $X_i$ and $Y_j$ is:

\[
  c[i, j] = \text{length of the LCS of } X_i \text{ and } Y_j
\]

with the recurrence:

\[
  c[i, j] =
  \begin{cases}
    0 & \text{if } i = 0 \text{ or } j = 0, \\
    c[i-1, j-1] + 1 & \text{if } i, j > 0 \text{ and } x_i = y_j, \\
    \max\bigl(c[i-1, j],\; c[i, j-1]\bigr) & \text{if } i, j > 0 \text{ and } x_i \ne y_j.
  \end{cases}
\]

We want to compute $c[m, n]$.

\begin{remark}{Overlapping Subproblems}{lcsov}
A naive recursion would solve the same subproblems an exponential number of times.
The bottom-up approach fills the table $c$ only once, by processing the
subproblems by increasing size (rows $i$ from $0$ to $m$, columns $j$ from $0$ to $n$).
\end{remark}

% ─────────────────────────────────────────────────────────────────────────────
\subsection{Bottom-up Algorithm and Complete Example}

For $X = \langle B,A,B,D,B,A \rangle$ and $Y = \langle D,A,C,B,C,B,A \rangle$
($m = 6$, $n = 7$), here is the complete table $c[i,j]$:

\begin{center}
\renewcommand{\arraystretch}{1.4}
\begin{tabular}{|c|c|c|c|c|c|c|c|c|c|}
& & \multicolumn{7}{c}{$j$} \\
$i$ & & \textbf{0} & \textbf{1 (D)} & \textbf{2 (A)} & \textbf{3 (C)} &
          \textbf{4 (B)} & \textbf{5 (C)} & \textbf{6 (B)} & \textbf{7 (A)} \\
\hline
\textbf{0}      & & 0 & 0 & 0 & 0 & 0 & 0 & 0 & 0 \\
\textbf{1 (B)}  & & 0 & 0 & 0 & 0 & 1 & 1 & 1 & 1 \\
\textbf{2 (A)}  & & 0 & 0 & 1 & 1 & 1 & 1 & 1 & 2 \\
\textbf{3 (B)}  & & 0 & 0 & 1 & 1 & 2 & 2 & 2 & 2 \\
\textbf{4 (D)}  & & 0 & 1 & 1 & 1 & 2 & 2 & 2 & 2 \\
\textbf{5 (B)}  & & 0 & 1 & 1 & 1 & 2 & 2 & 3 & 3 \\
\textbf{6 (A)}  & & 0 & 1 & 2 & 2 & 2 & 2 & 3 & \cellcolor{blue!20}\textbf{4} \\
\end{tabular}
\end{center}

The LCS has length \textbf{4} (value $c[6,7]$). By tracing back the decisions, we
obtain the LCS \textbf{ABBA}.

\begin{remark}{Reading the Table}{lcstable}
To fill $c[i,j]$: if $x_i = y_j$, come from the diagonal ($c[i-1,j-1]+1$).
Otherwise, take the maximum of the cell above ($c[i-1,j]$) and the cell
to the left ($c[i,j-1]$).
\end{remark}

% ─────────────────────────────────────────────────────────────────────────────
\subsection{Pseudocode}

\begin{algorithm}[H]
\caption{LCS-Length($X$, $Y$)}
\begin{algorithmic}[1]
\State $m \leftarrow X.\mathit{length}$;\; $n \leftarrow Y.\mathit{length}$
\State Create tables $c[0\ldots m,\; 0\ldots n]$ and $b[1\ldots m,\; 1\ldots n]$
\For{$i = 0$ \textbf{to} $m$}
    \State $c[i, 0] \leftarrow 0$
\EndFor
\For{$j = 0$ \textbf{to} $n$}
    \State $c[0, j] \leftarrow 0$
\EndFor
\For{$i = 1$ \textbf{to} $m$}
    \For{$j = 1$ \textbf{to} $n$}
        \If{$x_i = y_j$}
            \State $c[i, j] \leftarrow c[i-1, j-1] + 1$
            \State $b[i, j] \leftarrow \text{``}\nwarrow\text{''}$   \AlgComment{diagonal: common letter}
        \ElsIf{$c[i-1, j] \ge c[i, j-1]$}
            \State $c[i, j] \leftarrow c[i-1, j]$
            \State $b[i, j] \leftarrow \text{``}\uparrow\text{''}$   \AlgComment{up: advance in $X$}
        \Else
            \State $c[i, j] \leftarrow c[i, j-1]$
            \State $b[i, j] \leftarrow \text{``}\leftarrow\text{''}$ \AlgComment{left: advance in $Y$}
        \EndIf
    \EndFor
\EndFor
\State \Return $c$ and $b$
\end{algorithmic}
\end{algorithm}

\begin{complexity}{LCS-Length}{lcscplx}
\begin{itemize}[noitemsep]
  \item Table $c$ contains $(m+1)(n+1)$ entries, each calculated in $\Theta(1)$.
  \item \textbf{Time complexity: $\Theta(m \cdot n)$.}
  \item \textbf{Space complexity: $\Theta(m \cdot n)$.}
\end{itemize}
\end{complexity}

% ─────────────────────────────────────────────────────────────────────────────
\subsection{LCS Reconstruction}

Table $b$ stores the optimal choices: we trace it back recursively from
$b[m, n]$ toward the upper-left corner.

\begin{algorithm}[H]
\caption{Print-LCS($b$, $X$, $i$, $j$)}
\begin{algorithmic}[1]
\If{$i = 0$ \textbf{or} $j = 0$}
    \State \Return
\EndIf
\If{$b[i, j] = \text{``}\nwarrow\text{''}$}
    \State \Call{Print-LCS}{$b$, $X$, $i-1$, $j-1$}
    \State \textbf{print} $x_i$         \AlgComment{common letter: we include it}
\ElsIf{$b[i, j] = \text{``}\uparrow\text{''}$}
    \State \Call{Print-LCS}{$b$, $X$, $i-1$, $j$}
\Else
    \State \Call{Print-LCS}{$b$, $X$, $i$, $j-1$}
\EndIf
\end{algorithmic}
\end{algorithm}

\begin{complexity}{Print-LCS}{printlcscplx}
\begin{itemize}[noitemsep]
  \item Each recursive call decreases $i + j$ by at least 1.
  \item The recurrence satisfies $T(n) \le T(n-1) + \Theta(1)$, hence $T(n) = O(n)$.
  \item \textbf{Total complexity: $\Theta(|X| + |Y|)$}
        (the lower bound comes from $T(n) \ge T(n-2) + \Theta(1)$).
\end{itemize}
\end{complexity}

% ─────────────────────────────────────────────────────────────────────────────
\subsection{Backtracking Illustration}

For $X = \langle B,A,B,D,B,A \rangle$, $Y = \langle D,A,C,B,C,B,A \rangle$,
the path in table $b$ (illustrated below) identifies the letters \textbf{A, B, B, A}
= \textbf{ABBA} as the LCS:

\begin{center}
\begin{tikzpicture}[font=\small, >=Stealth,
    cell/.style={draw, minimum width=0.65cm, minimum height=0.55cm, inner sep=1pt}]

  % Column headers (Y)
  \foreach \c/\j in {{}/0, D/1, A/2, C/3, B/4, C/5, B/6, A/7}{
    \node[font=\bfseries\tiny] at (\j*0.7+0.35, 4.5) {\c};
  }
  % Row headers (X)
  \foreach \c/\i in {{}/0, B/1, A/2, B/3, D/4, B/5, A/6}{
    \node[font=\bfseries\tiny] at (-0.15, 4.0-\i*0.6) {\c};
  }

  % Fill the table with values and arrows
  % Row 0
  \foreach \j in {0,...,7}{
    \node[cell, fill=gray!10] at (\j*0.7+0.35, 4.0) {0};
  }
  % Row 1: 0 0 0 0 1 1 1 1
  \foreach \v/\j in {0/0, 0/1, 0/2, 0/3, 1/4, 1/5, 1/6, 1/7}{
    \pgfmathparse{int(\v)}
    \ifnum\pgfmathresult=1
      \node[cell, fill=green!15] at (\j*0.7+0.35, 3.4) {\v};
    \else
      \node[cell, fill=gray!10] at (\j*0.7+0.35, 3.4) {\v};
    \fi
  }
  % Row 2: 0 0 1 1 1 1 1 2
  \foreach \v/\j in {0/0, 0/1, 1/2, 1/3, 1/4, 1/5, 1/6, 2/7}{
    \pgfmathparse{int(\v)}
    \ifnum\pgfmathresult=2
      \node[cell, fill=keyblue!20] at (\j*0.7+0.35, 2.8) {\v};
    \else
      \node[cell, fill=gray!10] at (\j*0.7+0.35, 2.8) {\v};
    \fi
  }
  % Row 3: 0 0 1 1 2 2 2 2
  \foreach \v/\j in {0/0, 0/1, 1/2, 1/3, 2/4, 2/5, 2/6, 2/7}{
    \node[cell, fill=gray!10] at (\j*0.7+0.35, 2.2) {\v};
  }
  % Row 4: 0 1 1 1 2 2 2 2
  \foreach \v/\j in {0/0, 1/1, 1/2, 1/3, 2/4, 2/5, 2/6, 2/7}{
    \node[cell, fill=gray!10] at (\j*0.7+0.35, 1.6) {\v};
  }
  % Row 5: 0 1 1 1 2 2 3 3
  \foreach \v/\j in {0/0, 1/1, 1/2, 1/3, 2/4, 2/5, 3/6, 3/7}{
    \node[cell, fill=gray!10] at (\j*0.7+0.35, 1.0) {\v};
  }
  % Row 6: 0 1 2 2 2 2 3 4
  \foreach \v/\j in {0/0, 1/1, 2/2, 2/3, 2/4, 2/5, 3/6, 4/7}{
    \pgfmathparse{int(\v)}
    \ifnum\pgfmathresult=4
      \node[cell, fill=keygreen!25, font=\bfseries] at (\j*0.7+0.35, 0.4) {\v};
    \else
      \node[cell, fill=gray!10] at (\j*0.7+0.35, 0.4) {\v};
    \fi
  }

  % Traceback path arrows
  \draw[->, keyred, very thick, rounded corners=2pt]
    (7*0.7+0.35, 0.4)  % (6,7)=4, diagonal A match
    -- (6*0.7+0.35, 1.0)    % (5,6)=3, diagonal B match
    -- (4*0.7+0.35, 1.6)    % (4,4) up
    -- (4*0.7+0.35, 2.2)    % (3,4) diagonal B match
    -- (2*0.7+0.35, 2.8);   % (2,2) diagonal A match

  % Legend
  \node[keyred, font=\bfseries\small] at (3.5, -0.3)
        {LCS $=$ \textbf{ABBA} (length 4)};
\end{tikzpicture}
\end{center}

% ─────────────────────────────────────────────────────────────────────────────
\subsection{Summary: LCS}

\begin{tcolorbox}[colback=lightblue, colframe=keyblue, fonttitle=\bfseries,
                  title={Summary DP – Longest Common Subsequence}]
\begin{itemize}[noitemsep]
  \item \textbf{Choice}: if $x_i = y_j$, include this letter; otherwise advance in
        $X$ or in $Y$ (take the best of the two).
  \item \textbf{Optimal substructure}: 3 cases depending on whether $x_i = y_j$ or not
        (theorem above).
  \item \textbf{Recurrence}:
    \[
      c[i,j] =
      \begin{cases}
        0 & i = 0 \text{ or } j = 0,\\
        c[i-1,j-1]+1 & x_i = y_j,\\
        \max(c[i-1,j],\,c[i,j-1]) & x_i \ne y_j.
      \end{cases}
    \]
  \item \textbf{Complexity}: $\Theta(mn)$ in time and space.
  \item \textbf{Reconstruction}: table $b$ + \textsc{Print-LCS} in $\Theta(|X|+|Y|)$.
\end{itemize}
\end{tcolorbox}

% ─────────────────────────────────────────────────────────────────────────────
\subsection{General Overview of Dynamic Programming}

\begin{tcolorbox}[colback=lightyellow, colframe=orange!70!black, fonttitle=\bfseries,
                  title={DP Recipe in Three Steps}]
\begin{enumerate}
  \item \textbf{Identify the choices and the optimal substructure.}
  \item \textbf{Write the optimal value as a recurrence on subproblems.}
  \item \textbf{Solve efficiently}: top-down with memoization or bottom-up
        (without recalculating the same subproblems).
\end{enumerate}
\end{tcolorbox}

\begin{center}
\renewcommand{\arraystretch}{1.45}
\begin{tabular}{llcc}
\toprule
\textbf{Problem} & \textbf{Choice} & \textbf{Complexity} & \textbf{Table} \\
\midrule
Rod Cutting         & Left cut position               & $\Theta(n^2)$  & $n+1$ \\
Coin Change         & Denomination used               & $O(nW)$        & $W+1$ \\
Matrix Chains       & Parenthesis position            & $\Theta(n^3)$  & $n^2$ \\
LCS                 & Include $x_i = y_j$ / advance   & $\Theta(mn)$   & $(m+1)(n+1)$ \\
\bottomrule
\end{tabular}
\end{center}